% ---------------------------------------------------------------
% Shared style for "Vibe Coding in Practice" / "Vibe Coding 实战"
% Compile with XeLaTeX (latexmk -xelatex main.tex)
% ---------------------------------------------------------------
\usepackage[a4paper,top=24mm,bottom=26mm,left=24mm,right=24mm,headheight=14pt]{geometry}
\usepackage{fontspec}
\IfFontExistsTF{Charter}{\setmainfont{Charter}}{\setmainfont{TeX Gyre Schola}}
\IfFontExistsTF{Avenir Next}{
\setsansfont{Avenir Next}[UprightFont=* Regular,BoldFont=* Demi Bold,ItalicFont=* Italic,BoldItalicFont=* Demi Bold Italic]
\newfontfamily\displayfont{Avenir Next}[UprightFont=* Medium,BoldFont=* Demi Bold,Ligatures=TeX]
}{
\setsansfont{TeX Gyre Heros}
\newfontfamily\displayfont{TeX Gyre Heros}[Ligatures=TeX]
}
\IfFontExistsTF{Menlo}{\setmonofont{Menlo}[Scale=.82]}{\setmonofont{DejaVu Sans Mono}[Scale=.82]}

\usepackage{xcolor}
\usepackage{colortbl}
\definecolor{paper}{HTML}{F3EDE3}
\definecolor{ink}{HTML}{1C1B19}
\definecolor{accent}{HTML}{B4532A}
\definecolor{accentsoft}{HTML}{E9D3C4}
\definecolor{muted}{HTML}{6E675F}
\definecolor{rule}{HTML}{CBC1B2}
\definecolor{card}{HTML}{EAE2D5}
\definecolor{cardline}{HTML}{D6CBBB}
\definecolor{olive}{HTML}{5E6B4E}
\definecolor{ochre}{HTML}{B08A3E}
\pagecolor{paper}
\color{ink}

\usepackage{graphicx}
\usepackage{microtype}
\usepackage{needspace}
\usepackage{booktabs,tabularx,array,multirow,xltabular}
\usepackage{enumitem}
\setlist{itemsep=2pt,topsep=4pt,parsep=0pt,leftmargin=*}
\setlist[itemize]{label={\color{accent}\textbullet}}
\usepackage{amsmath,amssymb}
\usepackage{titlesec}
\usepackage{fancyhdr}
\usepackage{tikz}
\usetikzlibrary{arrows.meta,positioning,calc,shapes.geometric,shapes.misc,fit,backgrounds,decorations.pathreplacing}
\usepackage{pgfplots}
\pgfplotsset{compat=1.18}
\usepackage[most]{tcolorbox}
\usepackage{listings}
\usepackage{xurl}
\usepackage[hidelinks]{hyperref}
\hypersetup{colorlinks=true,linkcolor=accent,urlcolor=accent,citecolor=accent,
  pdftitle={Vibe Coding in Practice},pdfauthor={Bingjie Zhu}}

\setlength{\parindent}{0pt}
\setlength{\parskip}{6pt}
\linespread{1.08}

% ---------- Headings ----------
\titleformat{\section}[display]
  {\sffamily\color{ink}}
  {\color{accent}\small\bfseries\sectionlabel}
  {2pt}
  {\displayfont\fontsize{23}{27}\selectfont\raggedright\hyphenpenalty=10000\exhyphenpenalty=10000}
  [\vspace{3pt}{\color{rule}\titlerule[0.6pt]}]
\titlespacing*{\section}{0pt}{0pt}{8pt}
\titleformat{name=\section,numberless}[display]
  {\sffamily\color{ink}}{}{0pt}
  {\displayfont\fontsize{23}{27}\selectfont\raggedright\hyphenpenalty=10000\exhyphenpenalty=10000}
  [\vspace{3pt}{\color{rule}\titlerule[0.6pt]}]
\titleformat{\subsection}
  {\sffamily\bfseries\large\color{ink}}{\color{accent}\thesubsection}{0.7em}{}
\titlespacing*{\subsection}{0pt}{10pt}{3pt}
\titleformat{\subsubsection}
  {\sffamily\bfseries\normalsize\color{accent}}{}{0em}{}
\titlespacing*{\subsubsection}{0pt}{10pt}{2pt}
\renewcommand{\thesubsection}{\thesection.\arabic{subsection}}

% ---------- Header / footer ----------
\pagestyle{fancy}
\fancyhf{}
\renewcommand{\headrulewidth}{0pt}
\fancyhead[L]{\sffamily\scriptsize\color{muted}\MakeUppercase{Vibe Coding in Practice}}
\fancyhead[R]{\sffamily\scriptsize\color{muted}\leftmark}
\fancyfoot[C]{\sffamily\scriptsize\color{muted}\thepage}
\renewcommand{\sectionmark}[1]{\markboth{#1}{}}
\fancypagestyle{plain}{\fancyhf{}\fancyfoot[C]{\sffamily\scriptsize\color{muted}\thepage}}

% ---------- Small helpers ----------
\newcommand{\kicker}[1]{{\sffamily\small\bfseries\color{accent}\MakeUppercase{#1}}}
\newcommand{\lede}[1]{{\fontsize{12}{17.5}\selectfont #1\par}}
\newcommand{\term}[1]{\textbf{#1}}
\newcommand{\cmd}[1]{\texttt{\color{accent!80!ink}#1}}
\newcommand{\src}[1]{{\footnotesize\color{muted}[#1]}}

\newcommand{\moduleopen}[2]{%
  \noindent\begin{minipage}[c]{0.64\linewidth}\lede{#2}\end{minipage}\hfill
  \begin{minipage}[c]{0.32\linewidth}\includegraphics[width=\linewidth]{figures/#1}\end{minipage}\par\vspace{2pt}}
\newcommand{\keyideaname}{Note}
\newcommand{\logsname}{From my projects}
\newcommand{\modulename}{Chapter}
\newcommand{\sectionlabel}{\MakeUppercase{\modulename}~\thesection}
\tcbset{
  base/.style={enhanced,breakable,boxrule=0pt,arc=2pt,left=10pt,right=10pt,top=7pt,bottom=7pt,
    fonttitle=\sffamily\bfseries\small,coltitle=ink,before skip=8pt,after skip=10pt},
  tight/.style={before skip=4pt,after skip=6pt,top=5pt,bottom=5pt},
}
\newtcolorbox{keyidea}[1][\keyideaname]{base,unbreakable,colback=card,
  borderline west={2.4pt}{0pt}{accent},title={\color{accent}\MakeUppercase{#1}},
  attach title to upper,after title={\par\smallskip}}
\newtcolorbox{logs}[1][\logsname]{base,unbreakable,colback=accentsoft!45!paper,
  borderline west={2.4pt}{0pt}{ink},title={\MakeUppercase{#1}},attach title to upper,
  after title={\par\smallskip},fontupper=\small}
% Question / exercise box inside a chapter; the explanation follows as normal text.
\newtcolorbox{interact}[1]{base,unbreakable,colback=accentsoft!50!paper,colframe=accent,boxrule=0pt,
  borderline={0.9pt}{0pt}{accent,dash pattern=on 3pt off 2pt},
  title={\color{accent}#1},attach title to upper,after title={\par\smallskip}}
\newtcolorbox{promptbox}[1]{base,unbreakable,tight,colback=card,colframe=cardline,boxrule=0.5pt,
  title={\MakeUppercase{#1}},attach title to upper,after title={\par\smallskip},
  fontupper=\ttfamily\footnotesize}

% ---------- Code listings ----------
\lstdefinestyle{course}{
  basicstyle=\ttfamily\footnotesize\color{ink},
  backgroundcolor=\color{card},
  frame=single,framerule=0pt,framesep=6pt,xleftmargin=6pt,xrightmargin=6pt,
  commentstyle=\color{muted}\itshape,
  keywordstyle=\color{accent},
  stringstyle=\color{olive},
  showstringspaces=false,breaklines=true,columns=fullflexible,keepspaces=true,
  aboveskip=8pt,belowskip=8pt
}
\lstset{style=course}

% ---------- Tables ----------
\renewcommand{\arraystretch}{1.25}
\newcolumntype{L}[1]{>{\raggedright\arraybackslash}p{#1}}
\newcolumntype{Y}{>{\raggedright\arraybackslash}X}
\arrayrulecolor{rule}

% ---------- TikZ styles ----------
\tikzset{
  box/.style={draw=ink,line width=0.6pt,rounded corners=2pt,fill=paper,align=center,
    font=\sffamily\small,inner sep=5pt},
  abox/.style={box,draw=accent,fill=accentsoft!55!paper},
  soft/.style={box,draw=cardline,fill=card},
  arr/.style={-{Stealth[length=5pt]},line width=0.7pt,color=ink},
  aarr/.style={-{Stealth[length=5pt]},line width=0.9pt,color=accent},
  lbl/.style={font=\sffamily\scriptsize,color=muted,align=center},
}

\usepackage{xeCJK,fvextra,newunicodechar}
\setCJKmainfont{Noto Serif CJK SC}[BoldFont={Noto Serif CJK SC Bold},ItalicFont={Noto Sans CJK SC},BoldItalicFont={Noto Sans CJK SC Bold}]
\setCJKsansfont{Noto Sans CJK SC}[BoldFont={Noto Sans CJK SC Medium}]
\setCJKmonofont{Noto Sans CJK SC}[Scale=.92]
\setCJKfamilyfont{zhdisp}{Noto Sans CJK SC Medium}[BoldFont={Noto Sans CJK SC Bold}]
\let\displayfontlatin\displayfont
\renewcommand{\displayfont}{\displayfontlatin\CJKfamily{zhdisp}}
\linespread{1.2}\widowpenalty=10000\clubpenalty=10000
\setlength{\emergencystretch}{2em}
\newunicodechar{→}{\ensuremath{\rightarrow}}
\newunicodechar{↔}{\ensuremath{\leftrightarrow}}
\renewcommand{\sectionlabel}{第\,\thesection\,章}
\fancyhead[L]{\sffamily\scriptsize\color{muted}Vibe Coding 实战}
\hypersetup{pdftitle={Vibe Coding 实战},pdfauthor={朱炳杰},pdfsubject={以提示词组织任务定义、设计、实现与验证}}
\newcommand{\lesson}[2]{\clearpage\section{#2}\label{#1}\markboth{#2}{}}
\newcommand{\topic}[1]{\subsection*{#1}}
\newcommand{\diagram}[1]{\par\smallskip\includegraphics[width=\linewidth]{assets/diagrams/zh/#1.pdf}\par\smallskip}
\newtcolorbox{prompt}[1]{base,unbreakable,colback=card,colframe=cardline,boxrule=.5pt,
 title={#1},attach title to upper,after title={\par\smallskip},fontupper=\normalsize,fonttitle=\sffamily\bfseries\normalsize}
\newenvironment{concepts}{\par\smallskip\tabularx{\linewidth}{@{}L{29mm}Y@{}}\toprule 概念 & 定义\\\midrule}{\bottomrule\endtabularx\par}
\newcommand{\principle}[2]{\topic{#1}#2\par}
